\section{Routing, feedback loops, and termination}
\label{sec:control}

\subsection{Routing is code, not prose}
\label{sec:control-routing}

A conditional edge is where a run's next step is chosen, and the choice is a comparison written in
Python rather than a sentence written by a model. Both things are true at once, which is what makes
the arrangement worth describing carefully: the model's output decides which way the comparison
comes out, and the comparison decides where the run goes.

\begin{definition}[Conditional edge]
\label{def:conditional-edge}
A conditional edge is a function that receives the current state and returns the name of the next
node, or the terminal node \demoref{END}, which is ``a special node that represents a terminal
node''~\citep{langgraph-graph-api}. The function runs in the runtime's process and its return value
is the transition.
\end{definition}

Three responsibilities sit on three different components, and keeping them apart is the whole of
this section.

\bi

\item \textbf{The model proposes.} It emits a tool call, a plan, a candidate finding, or a judgement
about whether a candidate is acceptable. Every one of those is a value, not an action.

\item \textbf{The runtime records.} It merges the proposal into state through the reducers of
\cref{sec:state-reducers}, and it writes what happened to the trace.

\item \textbf{The router decides.} It reads named fields and returns a node name, applying the
transition rules and the limits we fixed when we wrote it.

\ei

Our loop's router is \demoref{after\_controller}, and what it reads is the point. It branches on
\demoref{status}, and then on whether the last message carried tool calls; there is no third
condition. It never reads the model's prose, so a model that writes ``I am confident the patch is
correct, submit it'' moves the run exactly as far as a model that writes nothing: the branch is
taken by the presence or absence of a tool call, not by the sentence. \demoref{after\_submit} is one
line, returning \demoref{END} when the submission was accepted and the controller otherwise.

The trace records who decided, and this is the field to read when a run surprises us. A
\demoref{routing} event carries \demoref{by="model"} when the model's output selected the branch
and \demoref{by="runtime"} when a limit did. \textbf{Neither value means the model chose the edge},
because the edge was chosen when we wrote the predicate; \demoref{by="model"} names the origin of
the value the predicate tested.

\subsection{The same graph, different paths}
\label{sec:control-paths}

Compilation fixes the graph and nothing else. \demoref{compile} adds three nodes, two conditional
edge sets and one unconditional edge, and the result is the same object on every run. What differs
between two runs on the same alert is the sequence of states, and therefore the sequence of edges
that the routers select.

\begin{figure}[htbp]
  \centering
  \begin{adjustbox}{width=\linewidth}% V4: the complete serial graph. Plan, execute, validate; on rejection, reflect, retain guidance,
% reset the attempt, and plan again. Two rows so the drawing stays inside the text width.
\begin{tikzpicture}[node distance=15mm and 7mm]
  \node[tnode] (in) {Input};
  \node[mnode, right=of in] (plan) {Plan agent};
  \node[mnode, right=of plan] (exec) {Reactive\\executor};
  \node[mnode, right=of exec] (val) {Validator};
  \node[tnode, right=20mm of val] (out) {Output};
  \node[rnode, below=of plan] (reset) {Retain guidance;\\reset attempt};
  \node[mnode, below=of val] (refl) {Reflection\\node};
  \node[tnode, align=center, below=of out, xshift=11mm] (stop) {Stop with\\unresolved result};
  \draw[flow] (in) -- (plan);
  \draw[flow] (plan) -- (exec);
  \draw[flow] (exec) -- (val);
  \draw[flow] (val) -- node[lab, above] {Accept} (out);
  \draw[back] (exec.north) to[out=120, in=60] node[lab, above] {planned step failed\\after retry} (plan.north);
  \draw[flow] (val) -- node[lab, left] {Reject; another\\attempt allowed} (refl);
  \draw[flow] (refl) -- (reset);
  \draw[back] (reset) -- (plan);
  \draw[flow] (val.south east) -- node[lab, pos=0.55, above, align=center] {Reject; attempt\\limit reached} (stop.west);
\end{tikzpicture}
\end{adjustbox}
  \caption{The complete serial graph, with both back edges. The upper edge revises a plan inside an
  attempt; the lower path runs the reflection step and begins a new attempt. Every edge in the
  drawing is a line of wiring code, and every label on one is a predicate over state.}
  \label{fig:control-graph}
\end{figure}

The unconditional edge is the one that made the loop a loop: the return from \demoref{tools} to
\demoref{controller}. Delete that line and the same three nodes become a chain. The conditional
edges are what give one compiled graph its different shapes at run time, and the run in
\cref{fig:seq-loop-run} traverses the controller four times without the graph changing once.

\subsection{Two returns to the planner}
\label{sec:control-backedges}

\Cref{fig:control-graph} has two back edges, and they answer different questions. One asks what to
do about a step that did not work; the other asks what to do about an attempt that did not work.
Collapsing them produces a system that either rewrites its plan over every transient tool error or
retries a whole investigation over one failed edit.

\begin{table}[htbp]
  \centering
  \begin{adjustbox}{width=\linewidth}\begin{tabularx}{\linewidth}{@{}L{0.155\linewidth}L{0.235\linewidth}XX@{}}
    \toprule
    \textbf{Back edge} & \textbf{Trigger} & \textbf{Retained} & \textbf{Reset} \\
    \midrule
    In-run plan revision & A planned step is still unsuccessful after the runtime's one free retry.
      The planner receives the failed subgoal and the error. & The attempt's observations, its
      progress so far, and the accounting. & The remaining subgoals, which the planner rewrites. \\
    New attempt after reflection & The validator rejects the candidate and the attempt limit
      permits another try. & Reflections, identity, the counters and the usage ledger. &
      Observations, plan, candidate patch and validation, all set explicitly. \\
    \bottomrule
  \end{tabularx}\end{adjustbox}
  \caption{The two back edges of \cref{fig:control-graph}. The retry that precedes the upper edge
  costs no model call, which is what makes it the runtime's decision rather than the model's.}
  \label{tab:control-backedges}
\end{table}

Ownership divides along the same line. The free retry is ours: the runtime reissues a failed tool
call once without consulting anyone, because a transient failure is not evidence about the plan.
The revision that follows a second failure is the model's proposal, recorded by the runtime and
merged into \demoref{plan}. The new attempt is ours again, because it is a comparison on the attempt
counter, and \textbf{a limit that the model can talk its way past is not a limit}.

\subsection{What actually stops a loop}
\label{sec:control-bounds}

Every loop in \cref{fig:control-graph} has an exit predicate, and no predicate is a guarantee. The
validator's acceptance condition says when an attempt may end; it says nothing about whether any
attempt will satisfy it, and a run whose candidate is rejected every time will keep going until
something other than acceptance stops it. Three limits do that work, and all three are comparisons
the runtime performs on fields it owns.

\bi

\item \textbf{The model-call budget.} Forty calls per run, and the controller checks it before it
issues a call rather than after. A limit checked afterwards is a post-mortem: the call has already
been paid for. Exceeding it sets \demoref{budget\_exhausted} and the router sends the run to
\demoref{END}.

\item \textbf{The stall check.} Three consecutive observations with identical results end the run
with \demoref{no\_progress}. Comparing results rather than requests is deliberate, because it also
catches six distinct searches that all return nothing, which is the repeated-action failure mode
that ReAct reports for its own trajectories \citep[\S3.3]{yao2023react}.

\item \textbf{The attempt limit.} Two attempts, counted in a field that survives the reflection
reset. An attempt counter that reset with the attempt would bound nothing.

\ei

The framework contributes one more, underneath all of these: a recursion limit on the number of
super-steps a single execution may run, which raises an error when it is
reached~\citep{langgraph-graph-api}. It is a failsafe against a wiring mistake and not a control on
the investigation, because it knows nothing about alerts, attempts or costs. We prefer our limits to
fire first, and the budget of forty is set low enough that they do.

\subsection{What a trace establishes}
\label{sec:control-termination}

A run ends in one of five ways, and the vocabulary is closed on purpose: \demoref{accepted},
\demoref{unresolved}, \demoref{budget\_exhausted}, \demoref{no\_progress} and \demoref{error}. The
trace writer refuses a status outside that set, and it also refuses the opposite mistake. Closing a
trace that never recorded a terminal event raises \demoref{NoTerminalStatus}, because
\textbf{a run that never said how it ended is a defect in our instrumentation rather than an
outcome we can report}.

That is what we control. What a trace lets a reader conclude is a separate question, and it has
three answers rather than five.

\begin{table}[htbp]
  \centering
  \begin{adjustbox}{width=\linewidth}\begin{tabularx}{\linewidth}{@{}L{0.31\linewidth}X@{}}
    \toprule
    \textbf{What the trace records} & \textbf{What it supports} \\
    \midrule
    A verdict, then a terminal event with status \demoref{accepted} & A completed assessment under
    the runtime's acceptance rules, whose quality is a separate question from its completion. \\
    A terminal event with a stop status and its reason & A declared end to execution without a
    completed assessment, with whatever partial evidence the run had gathered. \\
    No terminal event & Completion has not been established by this record. The run may still be
    active, may have been interrupted, or may have written an incomplete trace. \\
    \bottomrule
  \end{tabularx}\end{adjustbox}
  \caption{Three readings of a trace. The third is a statement about the evidence and not about the
  run, and treating it as \demoref{error} attributes a cause the record does not carry.}
  \label{tab:control-outcomes}
\end{table}

Three distinctions keep those readings honest. Errors and outcomes are not the same thing: a
refused tool call is an observation inside a run that can still end in \demoref{accepted}, and a
process killed between two events leaves a run with no status at all. The finding and the execution
status are not the same thing either, because an assessment that reports insufficient information is
a completed assessment, while a declared stop can carry partial evidence and no finding. And
reaching \demoref{END} establishes only that the graph stopped: the reason lives in
\demoref{status}, which we set.

Termination therefore stays where routing is. A model's statement that it is finished is a proposal,
like every other value it produces; the controller applies the completion rules, sets the status,
and the trace records what was set. Stopping conditions and token-budget accounting are taken up
in their own right in the coming weeks.

\subsection{The demo: three endings}
\label{sec:control-demo}

The recorded run of \cref{fig:seq-loop-run} ends the first way. Its seventeenth event is a
\demoref{terminal} record with status \demoref{accepted} and the run's usage totals, written after
\demoref{submit} reported that all three mechanical checks passed, and the routing event just before
it carries \demoref{by="runtime"}. Nothing in that trace is a model's opinion about whether the work
was done.

A budget-exhausted run ends the second way, and its shape differs in one instructive place. The
controller's last event is not a model call, because the check that fired precedes the call: the
trace shows a \demoref{routing} event with decision \demoref{budget\_exhausted} and
\demoref{by="runtime"}, then the terminal event, and no fifth model call between them. The evidence
the run had gathered is still in \demoref{observations}, and the absence of a finding is not the
absence of work.

The third ending is the one we cannot label. Interrupt a run between its tool result and its next
model call and the trace stops there, with events that are all individually valid and no
\demoref{terminal} record at the end. What we may say about that file is that no ending was
observed. What we may not say is which ending would have arrived, and an excerpt used in teaching
is worth marking with how it was produced.

\subsection{Exercises}

\begin{enumerate}
\item \demoref{after\_controller} reads two fields. Write a model reply that would make a reader
  expect the run to end and that in fact routes to \demoref{tools}, and say which field decided.
\item The budget check sits before the model call. Move it after, and describe the two ways the
  resulting trace would differ from the current one on a run that exhausts the budget.
\item A colleague proposes ending a run when the model writes \verb|DONE| in its message. Give one
  run in which that predicate ends the run too early and one in which it never fires, and name the
  field each failure would have read instead.
\item Take a trace whose last event is a tool result. List everything you can state about how that
  run ended, and mark each statement with the event that supports it.
\end{enumerate}

\subsection{Reading}

\citet{langgraph-graph-api}, the sections on conditional edges, \demoref{END} and the
recursion limit. \S3.3 of \citet{yao2023react} supplies the repeated-action failure mode that the
stall check is built against, with its own trajectory counts.

For reference while writing HW1: the terminal statuses and trace events of \demoref{trace.py}, which
are the vocabulary every claim in \cref{tab:control-outcomes} is made in.

\paragraph{Looking ahead.} Routing decides which node runs next, and one of those nodes executes
what the model asked for. The last section of these notes takes that step apart. We ask what a tool
contract contains beyond the schema the model sees, what the dispatcher does between a request and a
result, which artifact proves that a function was entered, and which component refuses a request
that the run is not allowed to make.
